\documentclass{article}
\usepackage[T1]{fontenc}
\usepackage{lmodern}
\usepackage{graphicx}
\usepackage{amsmath,amssymb}
\usepackage[a4paper,margin=2cm]{geometry}
\usepackage[absolute,overlay]{textpos}
\setlength{\TPHorizModule}{1mm}
\setlength{\TPVertModule}{1mm}
\usepackage[indonesian]{babel}
\usepackage{hyperref}
\setlength{\emergencystretch}{2em}
% unit: Halaman 3

\begin{document}

% === Halaman 3 (llm) ===

\section*{Serangan (\textit{attack})}

\begin{itemize}
    \item \textbf{Serangan} diartikan sebagai setiap usaha (\textit{attempt}) atau percobaan yang dilakukan oleh pihak lawan untuk menemukan kunci atau menemukan plainteks dari cipherteksnya.
    \item Asumsi: penyerang mengetahui algoritma kriptografi yang digunakan
\end{itemize}

\textbf{Prinsip Kerckhoff}: Semua algoritma kriptografi harus publik; hanya kunci yang rahasia.

Jadi, keamanan sistem kriptografi terletak pada \textcolor{red}{kunci}!

\end{document}
